\scanpage{69}
\mspage{- 30 -}

D'autre part, on a, en tenant compte de (26)

\[
\begin{aligned}
\frac{d}{dt}(tU_{kj}(\handwritten{x}_0+t(\handwritten{x}-\handwritten{x}_0);\handwritten{v}(t)))&=U_{kj}(\handwritten{x}_0+t(\handwritten{x}-\handwritten{x}_0);\handwritten{v}(t))+\\
&\quad+t\handwritten{\sum_i}(x_i-x_{i0})\frac{\handwritten{\partial}U_{kj}}{\handwritten{\partial} x_i}(\handwritten{x}_0+t(\handwritten{x}-\handwritten{x}_0);\handwritten{v}(t))\\
&\quad+t\handwritten{\sum_{i,h}}(x_i-x_{i0})\frac{\handwritten{\partial} U_{kj}}{\handwritten{\partial} y_h}(\handwritten{x}_0+t(\handwritten{x}-\handwritten{x}_0);\handwritten{v}(t))U_{ih}(\handwritten{x}_0+t(\handwritten{x}-\handwritten{x}_0);\handwritten{v}(t))
\end{aligned}
\]

d'où, en faisant intervenir (24), on voit que les fonctions

\[
w_{kj}(t)=\frac{\handwritten{\partial} v_j}{\handwritten{\partial} x_k}-tU_{kj}(\handwritten{x}_0+t(\handwritten{x}-\handwritten{x}_0);\handwritten{v}(t))\qquad
(1\leq j\leq q,\ 1\leq k\leq p)
\]

vérifient le système linéaire et homogène des $pq$ équations différentielles

\[
\tag*{(27)}\frac{dw_{kj}}{dt}=\handwritten{\sum_{i,h}}(x_i-x_{i0})
\frac{\handwritten{\partial} U_{ij}}{\handwritten{\partial} y_h}w_{kh}
\]

Or $w_{kj}(0)=0$, car pour $t=0$, $\handwritten{v}(0;\handwritten{x}_0,\handwritten{x},\handwritten{y}_0)=\handwritten{y}_0$, donc $\frac{\handwritten{\partial} v_j}{\handwritten{\partial} x_k}(0)=0$ ; comme $0$ est la seule solution de (27) qui s'annule pour $t=0$, on a identiquement $w_{kj}(t)=0$ pour $|t|<2$, et en particulier pour $t=1$ ; mais si on pose $\handwritten{u}(\handwritten{x},\handwritten{x}_0,\handwritten{y}_0)=\handwritten{v}(1;\handwritten{x}_0,\handwritten{x},\handwritten{y}_0)$, cela signifie que l'on a

\[
\frac{\handwritten{\partial}u_j}{\handwritten{\partial}x_i}=U_{ij}(\handwritten{x};\handwritten{u}(\handwritten{x}))\qquad
(1\leq i\leq p,\ 1\leq j\leq q)
\]

pour $\handwritten{x}\handwritten{\in}U$. De ces relations on déduit que si $\handwritten{z}(t)=\handwritten{u}(\handwritten{x}_0+t(\handwritten{x}-\handwritten{x}_0))$, on a

\[
\handwritten{z}'(t)=\underline{U}(\handwritten{x}_0+t(\handwritten{x}-\handwritten{x}_0);\handwritten{z}(t)).(\handwritten{x}-\handwritten{x}_0)
\]

pour $|t|<2$, et $\handwritten{z}(1)=\handwritten{v}(1)$ ; donc $\handwritten{z}(t)=\handwritten{v}(t)$\unsure{.|,} pour $|t|\leq 1$, et en particulier $\handwritten{u}(\handwritten{x}_0)=\handwritten{z}(0)=\handwritten{v}(0)=\handwritten{y}_0$ ; ce qui achève la démonstration.

On pourrait évidemment ici passer de la forme non homogène à la forme homogène du th. de Frobenius, avec l'algèbre extérieure, parler de la primitive (locale) d'une forme extérieure holomorphe, etc..; est-ce le lieu, ou,vaut-il mieux attendre les variétés analytiques pour le faire ?
% archive sheet number: 68

